\lesson{5}{Sep 27 2021 Mon (09:06:22)}{Plot, Pacing, and Point of View}{Unit 1}

Sometimes, our interpretations of a person or a situation are incorrect because
we do not have all of the information, we see what we want to see, or we
interpret based on our own experiences. Whatever the reasons, at some point,
our expectations or interpretations often come crashing down, and we must face
reality.

\subsubsection*{Elements of the Short Story}

A short story is simply shorter and must do the work of a longer story much
more quickly. Take a look at the list below to review the elements of a short
story.

\begin{definition}[Plot]
    The story line or action and events in a story.
\end{definition}

\begin{definition}[Conflict]
    The main struggle or the many struggles a character faces; without
    conflict, there is no story.
\end{definition}

\begin{definition}[Setting]
    The location and time in which a story takes place. Sometimes, it is simply
    the backdrop of the story, but sometimes it functions as a
    character, influencing the meaning and mood of the story.
\end{definition}

\begin{definition}[Character]
    The individuals in a work of fiction. Characters can be well developed, or
    they can be minor characters about whom we only know one or
    two characteristics.
\end{definition}

\begin{definition}[Point of View]
    The angle or perspective from which the narrator tells the story. The point
    of view affects what a reader knows about the characters and
    action of the story.
\end{definition}

\begin{definition}[Theme]
    The meaning of the work or the life lesson that readers should learn as a
    result of reading the story.
\end{definition}

The plot is what happens in the story. This element of fiction can be broken
down into five parts:

\begin{itemize}
    \item \textbf{Exposition:} The introduction to characters, setting, and
        situation of the story.
    \item \textbf{Rising Action:} The development of complications and
        conflict. This part of the plot diagram may be the longest side of the
        triangle if the rising action is the bulk of the story. Or it may be
        shorter if the climax occurs early in the story.
    \item \textbf{Climax or Crisis:} The moment of greatest emotional intensity
        or tension that usually comes at the turning point in a story.
    \item \textbf{Falling Action:} The action that follows the climax.
    \item \textbf{Resolution:} The revelation of the outcome of the conflict,
        also called the \textbf{"denouement"} which means \textbf{"tying up"}.
\end{itemize}

\subsubsection*{Different Types of Views}

\begin{definition}[First Person]
    Narrator is a character and participates in the action.
    
    \paragraph{Effect} Narrator and, therefore, reader have limited knowledge
    of events or thoughts and feelings of other characters, but complete
    knowledge of the narrator’s experience. Narrator confides in reader; we
    often feel close to a first person narrator. This point of view uses
    pronouns like \textbf{"I"}, \textbf{"we"}, and \textbf{"me}" because the
    narrator is part of the action.
\end{definition}

\begin{definition}[Objective]
    Narrator is unidentified, a detached observer.
    
    \paragraph{Effect} Like a video camera that follows the characters and the
    action, this type of narrator reports on events and lets the reader supply
    the meaning. We have to draw our own conclusions about what characters are
    thinking or feeling.
\end{definition}

\begin{definition}[Omniscient]
    Narrator is knows all.
    
    \paragraph{Effect} The narrator knows what everyone is thinking and feeling
    and reveals that information to the reader. An omniscient narrator may
    provide his or her opinion of events and characters, or he or she may let
    the reader come to his or her own conclusions. Thanks to the insights of a
    third-person omniscient narrator, we might feel sympathy or anger toward
    another character because we know the other "pieces" of the story. This
    point of view uses pronouns like \textbf{"he"}, \textbf{"she"}, and
    \textbf{"they"} because the narrator is not part of the action.
\end{definition}

\begin{definition}[Limited omniscient]
    Narrator knows the thoughts, feelings, and experiences of a single character.
    
    \paragraph{Effect} The narrator shares his or her limited knowledge with
    the reader, leaving some things unknown. This point of view may
    use pronouns like \textbf{"he"}, \textbf{"she"} and \textbf{"they"} because
    the narrator is not part of the action.
\end{definition}

\newpage
